\documentclass[UTF8,a4paper,12pt]{ctexart}

\usepackage[margin=1in]{geometry}
\usepackage{array}
\usepackage{booktabs}
\usepackage{float}
\usepackage{xcolor}
\usepackage{hyperref}
\setlength{\emergencystretch}{3em}

\definecolor{reportlink}{RGB}{65, 145, 210}

\hypersetup{
  colorlinks=true,
  linkcolor=reportlink,
  urlcolor=reportlink,
  citecolor=reportlink
}

\newcommand{\project}{DiPECS}

\title{\project{} 可行性报告}
\author{\project{} Team}
\date{2026 年 4 月 15 日}

\begin{document}

\maketitle

\begin{abstract}
\project{} 是一个面向 Android/Linux 端侧环境的意图预测与动作治理原型。项目希望在不修改普通用户
操作系统内核、不让模型直接控制设备的前提下，完成一条“采集、脱敏、聚合、决策、策略审查、动作执行、
审计回放”的闭环。本报告从数据采集、系统架构、决策模型、动作执行、隐私约束和实验评估等方面分析
该项目的技术可行性。结论是：基于 Android 公开 API、Rust 核心链路、离线 replay 和受控动作边界，
\project{} 可以先实现一个可运行、可验证的端侧 AIOS 原型；高权限采集、真实系统预加载和复杂语义模型
应作为后续阶段逐步推进。
\end{abstract}

\newpage
{\hypersetup{linkcolor=black}\tableofcontents}
\newpage

\section{项目概述}

\project{} 的目标不是实现一个通用聊天助手，而是在 Android/Linux 场景下验证端侧 AIOS 的核心机制：
系统如何从多源事件中获得上下文，如何在本地完成隐私脱敏，如何把上下文交给本地规则或云端模型进行
意图判断，以及如何确保模型建议不会绕过本地策略直接执行。AIOS 相关工作已经说明，LLM agent 需要调度、
上下文、存储和访问控制等系统级支撑；\project{} 则把这个问题收敛到 Android/Linux 端侧场景 \cite{aios}。

当前项目采用 Kotlin + Rust 的工程路线。Kotlin 负责 Android 应用层能力，包括权限申请、通知监听、
前台服务、设备状态采集和动作桥接；Rust 负责统一事件模型、隐私空气隙、窗口聚合、决策路由、策略审查、
动作生命周期、离线 replay 和审计记录。两侧通过 JSONL/schema 边界传输结构化数据，使同一套 Rust
核心逻辑既能处理真实 Android trace，也能处理离线测试 fixture。

系统主链路可以概括为：

\begin{center}
\small
RawEvent $\rightarrow$ PrivacyAirGap $\rightarrow$ SanitizedEvent $\rightarrow$
StructuredContext $\rightarrow$ DecisionRouter $\rightarrow$ PolicyEngine $\rightarrow$
ActionLifecycle $\rightarrow$ AuditRecord
\end{center}

这条链路中的每一步都有明确输入输出，因此项目可以通过 staged development 方式逐步实现：先跑通离线
replay 和本地规则，再接入 Android 公共 API trace，最后补充云端模型、Android bridge 和更复杂的预测后端。

\section{系统设计可行性}

\subsection{总体架构}

\project{} 的架构可以拆为六层：

\begin{table}[H]
\centering
\small
\begin{tabular}{p{0.22\linewidth}p{0.28\linewidth}p{0.40\linewidth}}
\toprule
层级 & 代表模块 & 主要职责 \\
\midrule
采集层 & Android collector, aios-collector & 采集应用切换、通知、屏幕、电量、进程资源和 JSONL trace。 \\
隐私层 & PrivacyAirGap & 将 raw event 转为 sanitized event，丢弃通知正文、文件路径等敏感信息。 \\
上下文层 & WindowAggregator & 按时间窗口构造 StructuredContext 和 ContextSummary。 \\
决策层 & DecisionRouter, RuleBased, CloudLlm & 根据复杂度、隐私分数和后端状态选择本地或云端决策。 \\
治理层 & PolicyEngine, ActionLifecycle & 检查 capability、风险、置信度、target 和 action schema。 \\
执行/验证层 & aios-action, OfflineAdapter, replay & 执行低风险动作或离线模拟，生成 audit record 和 hash。 \\
\bottomrule
\end{tabular}
\caption{\project{} 的分层设计}
\end{table}

这种分层的可行性来自两个事实。第一，Android 和 Linux 已经能提供足够的低粒度事件，支持最小原型；
第二，Rust 中的核心模块不依赖 Android UI，可以用离线 trace 进行确定性测试。即使 Android 真机能力
尚未完全接入，系统仍可通过 replay 验证隐私、路由、策略和审计链路。

\subsection{数据采集可行性}

Android 普通应用无法任意读取其他应用内部状态，但仍可在用户授权下使用公开 API 获取有限而稳定的上下文。
UsageStats 可以提供应用前后台切换趋势；NotificationListener 可以获取通知来源、类别和出现时间；
BatteryManager、ConnectivityManager 和 PowerManager 可以提供电量、网络和屏幕状态 \cite{android-usagestats,android-notification,android-battery}。
这些信号虽然不能完整
表达用户意图，但足以构造第一阶段的行为窗口。

更高权限的数据源，如 Accessibility、ADB/Shizuku、Binder trace、eBPF 或 root 环境，可以提高观测精度，
但部署成本和隐私风险显著更高。因此本项目的可行路径是：默认主链路依赖公开 API 和本机 Linux 采集；
高权限数据源只作为实验增强，不作为最小闭环的前提 \cite{android-accessibility,aosp-ebpf}。

\subsection{隐私脱敏可行性}

端侧上下文天然包含敏感信息。通知标题和正文可能包含联系人、验证码或工作内容；文件路径可能包含项目名、
客户名或个人目录；精确地点可能直接暴露用户轨迹。因此系统不能把 raw event 直接交给云端模型。

\project{} 采用本地确定性脱敏。RawEvent 只存在于 collector 到 core 的短路径中，经过 PrivacyAirGap 后，
后续模块只处理 SanitizedEvent。例如，通知正文被转换为文本长度、书写系统和 semantic hint；文件路径只用于
扩展名类别判断；notification key 等高风险字段被丢弃。这种设计牺牲一部分语义完整性，但换来清晰的隐私边界
和可测试的泄漏防护。

\subsection{上下文聚合可行性}

单个事件很难支撑意图判断。用户是否即将打开某个应用，通常要结合通知、前台应用、屏幕状态、电量和最近进程状态
共同判断。因此项目使用时间窗口聚合，把 SanitizedEvent 序列转为 StructuredContext。

窗口聚合的计算成本低，适合端侧运行；输出结构稳定，适合 replay；summary 字段又能降低后端决策复杂度。
这使得 RuleBased、CloudLlm 和未来的统计预测模型都可以共享同一输入格式，避免为不同后端重复处理 raw event。

\section{决策模型可行性}

\subsection{RuleBased 基线}

RuleBased 是最可行的第一阶段后端。它不依赖网络、不需要训练数据，行为可解释，适合作为默认降级路径。
当前合理边界是：RuleBased 只生成低风险、能力范围内的动作，例如 NoOp、KeepAlive 和 ReleaseMemory。
它不应直接生成 PreWarmProcess 或 PrefetchFile，因为这些动作需要更明确的 target 语义和更高能力等级。

RuleBased 的价值不在于预测能力最强，而在于建立系统基线：给定同一 StructuredContext，它应产生稳定的
IntentBatch，并且不应持续触发 capability denial。这个基线为后续云端模型和统计模型提供对照。

\subsection{统计预测模型}

在 RuleBased 之后，最自然的增强是统计转移模型。移动行为预测研究表明，用户应用切换存在明显时序规律。
项目可以从 replay trace 或真实 Android trace 中统计 foreground app 之间的转移概率，形成
$P(a_{t+1} \mid a_t)$ 或 $P(a_{t+1} \mid a_{t-k}, \ldots, a_t)$。

统计模型实现成本低，推理开销小，且容易解释。它可以作为 LocalEvaluator 的第一版：在上下文中输出 top-k
候选 app 和置信度，再由 PolicyEngine 判断是否允许 KeepAlive 或更高级动作。若统计模型不能显著优于
RuleBased，复杂序列模型也不应过早投入。

\subsection{云端 LLM 后端}

云端 LLM 的可行性体现在语义理解和规则编排上。对于包含多个 semantic hint 的复杂上下文，LLM 可以基于
结构化 JSON 生成 IntentBatch，并给出风险等级、置信度和 rationale tags。当前设计中，CloudLlm 是可选后端：
只有配置完整且 router 判断上下文复杂度足够时才调用；失败后回退到 RuleBased，并由 circuit breaker 防止
连续错误放大。OS-Copilot/FRIDAY 这类系统表明，LLM 可以组织跨应用任务和工具调用；但在 \project{} 中，
这类能力必须被本地 policy 边界约束 \cite{os-copilot,os-copilot-page}。

这种设计降低了云端依赖风险。网络不可用、API 失败或模型输出格式错误都不会阻塞本地链路。但同时也说明，
云端模型不能作为唯一控制面；所有云端输出仍必须经过本地策略和动作生命周期。

\section{动作执行可行性}

动作执行是项目风险最高的部分。相比预测“用户可能打开某个应用”，执行“预热、保活、预取、释放内存”等动作
会直接影响系统资源。为了降低风险，\project{} 将动作分为建议、授权和执行三个阶段。

\begin{table}[H]
\centering
\small
\begin{tabular}{p{0.22\linewidth}p{0.34\linewidth}p{0.34\linewidth}}
\toprule
动作 & 当前可行语义 & 主要限制 \\
\midrule
NoOp & 不做动作，只记录系统观察结果 & 无直接收益，但适合安全降级。 \\
KeepAlive & 对已出现 app 或系统任务做轻量保活提示 & Android 普通应用能做的真实保活能力有限。 \\
ReleaseMemory & 释放系统或缓存类资源，降低压力 & package 级释放需与 Android bridge 能力对齐。 \\
PreWarmProcess & 预热应用或本系统资源 & 第三方 app 预热受平台限制，不能默认承诺。 \\
PrefetchFile & 对 url/uri 或可访问内容做预取 & 需要明确 target 权限和 policy 放行规则。 \\
\bottomrule
\end{tabular}
\caption{动作类型的可行性与限制}
\end{table}

因此，最可行的执行路线是先用 OfflineAdapter 验证 action lifecycle，再逐步接入 Android bridge。真实执行层应
优先支持 Android-safe 的动作，例如本应用资源预热、公开 content uri 预取、cache trim 或用户可见通知提示。
对第三方进程强制预热、系统级内存调度等能力，应明确标记为 root/AOSP 阶段目标。

\section{验证与评估可行性}

\project{} 的一个优势是可以离线 replay。Replay 不访问真实 Android、不访问网络，也不执行真实系统动作，而是
读取 JSONL trace，经过相同的 Rust pipeline，生成 intent、policy decision、execute record 和 audit hash。
这使项目可以在没有真机或云端的情况下持续验证核心逻辑。

移动端预加载系统的评估不能只看预测是否命中，还要看启动延迟、能耗和误预加载成本。FALCON 对 app prelaunch
的评估方式可以作为后续实验设计参考 \cite{falcon}。

评估指标应分为三组：

\begin{itemize}
  \item \textbf{预测指标}：top-1/top-k 命中率、置信度校准、不同后端相对 RuleBased 的提升。
  \item \textbf{系统指标}：动作授权数、拒绝数、noop rate、policy denial reason、执行失败数、replay hash 稳定性。
  \item \textbf{资源指标}：动作触发后的启动延迟变化、CPU/内存开销、误预测导致的资源浪费率。
\end{itemize}

第一阶段应优先使用系统指标和 replay 稳定性，因为它们直接验证系统边界是否正确。只有在动作桥接更完整后，
再重点评估启动延迟和资源收益。

\section{风险分析}

\subsection{权限与部署风险}

Android 公开 API 能力有限，高权限数据源部署复杂。应对策略是把公开 API 作为默认主路径，把 root、eBPF、
Binder trace 明确为实验增强。报告、代码和展示都不应把高权限能力写成已经完成的默认依赖。

\subsection{隐私与语义损失风险}

脱敏会丢失通知正文、文件路径和应用内文本，可能降低模型判断能力。应对策略是用 semantic hint、类别标签和
summary 字段补偿部分语义，同时用 replay/audit 确保脱敏边界不被绕过。

\subsection{云端不稳定风险}

云端 LLM 可能因为网络、配额或格式错误失败。应对策略是保持 RuleBased fallback、配置 circuit breaker，
并要求云端只输出结构化 JSON，不直接控制设备。

\subsection{动作收益不确定风险}

即使预测正确，动作也未必带来可感知收益。例如某些应用启动本来很快，预加载反而浪费资源。应对策略是引入
confidence threshold、预算限制和 wasted action rate，用实验数据决定哪些动作值得保留。

\section{阶段计划}

\begin{table}[H]
\centering
\small
\begin{tabular}{p{0.16\linewidth}p{0.34\linewidth}p{0.38\linewidth}}
\toprule
阶段 & 目标 & 可验证结果 \\
\midrule
Phase 1 & 打通离线 replay 和 RuleBased 基线 & 同一 trace 产生稳定 intent、policy 和 audit hash。 \\
Phase 2 & 接入 Android public-API trace & 真机或模拟器 JSONL 能进入 Rust pipeline。 \\
Phase 3 & 完善 cloud/rule/local 后端路由 & 云端失败可回退，privacy-sensitive context 不上传。 \\
Phase 4 & 增加统计预测后端 & 与 RuleBased 比较 top-k 命中率和 noop rate。 \\
Phase 5 & 强化 Android-safe action bridge & 真实动作有审计记录，并能量化收益与开销。 \\
\bottomrule
\end{tabular}
\caption{阶段性开发路线}
\end{table}

\section{结论}

综合分析，\project{} 的最小闭环具备可行性。Android 公开 API 足以提供第一阶段上下文；Rust 核心可以稳定承担
脱敏、聚合、策略和 replay；RuleBased 可以作为本地安全基线；CloudLlm 和统计模型可以作为后续增强；
ActionLifecycle 和 audit record 能限制动作风险并提供可验证记录。

项目当前需要避免过度承诺高权限采集和系统级资源调度。可行的发展策略是先保证数据链路、隐私边界和动作治理
正确，再逐步证明预测模型和真实动作是否带来性能收益。这样既能形成可演示原型，也能保留继续扩展到 root/AOSP
环境的空间。

\begin{thebibliography}{99}

\bibitem{android-usagestats}
Android Developers.
\newblock UsageStatsManager.
\newblock \href{https://developer.android.com/reference/android/app/usage/UsageStatsManager}{Android API reference}.

\bibitem{android-notification}
Android Developers.
\newblock NotificationListenerService.
\newblock \href{https://developer.android.com/reference/android/service/notification/NotificationListenerService}{Android API reference}.

\bibitem{android-battery}
Android Developers.
\newblock BatteryManager.
\newblock \href{https://developer.android.com/reference/android/os/BatteryManager}{Android API reference}.

\bibitem{android-accessibility}
Android Developers.
\newblock AccessibilityService.
\newblock \href{https://developer.android.com/reference/android/accessibilityservice/AccessibilityService}{Android API reference}.

\bibitem{aosp-ebpf}
Android Open Source Project.
\newblock Extend the kernel with eBPF.
\newblock \href{https://source.android.com/docs/core/architecture/kernel/bpf}{AOSP documentation}.

\bibitem{aios}
Kai Mei et al.
\newblock AIOS: LLM Agent Operating System.
\newblock \href{https://arxiv.org/abs/2403.16971}{arXiv:2403.16971}.

\bibitem{falcon}
Tingxin Yan et al.
\newblock FALCON: Fast App Launching via Predictive User Context.
\newblock \href{https://www.microsoft.com/en-us/research/publication/fast-app-launching-for-mobile-devices-using-predictive-user-context/}{Microsoft Research page}.

\bibitem{os-copilot}
Zhiyong Wu et al.
\newblock OS-Copilot: Towards Generalist Computer Agents with Self-Improvement.
\newblock \href{https://arxiv.org/abs/2402.07456}{arXiv:2402.07456}.

\bibitem{os-copilot-page}
OS-Copilot project page.
\newblock \href{https://os-copilot.github.io/}{os-copilot.github.io}.

\bibitem{dipecs-repo}
\project{} repository.
\newblock Rust crates for spec, core, collector, agent, action, CLI and daemon.

\end{thebibliography}

\end{document}
